% Ad Atticum 14.16 --- headnote
% ad-atticum-14-16, 3 May 44 BC, in Puteolano

\textit{Cicero to Atticus, written at Puteoli, in the Cluvian
gardens, on 3 May 44 BC --- Perseus dateline \textit{Scr. in
Puteolano sive hortis Cluvianis v Non. Mai. a. 710 (44)}. The
letter is dictated as Cicero boards a small rowing-boat at the
Cluvian estate, with a stopover planned at Paetus's table
(\textit{tyrotarichum} --- ``cheese-and-salt-fish,'' the comic
shorthand for a modest meal) and then on to Pompeii. The bay of
Naples is at once delicious and uninhabitable: \textit{o loca
ceteroqui valde expetenda, interpellantium autem multitudine
paene fugienda}.}

\textit{The substantive middle section returns to the praise of
Dolabella begun in 14.15, with two Greek tags piled up:
\textit{aristeian} (an Iliadic word --- a hero's individual feat
of arms) and \textit{anathe\=or\=esis} (the ``reconsideration''
of 14.15, repeated). Brutus, Cicero half-jokes, could now wear a
gold crown through the Forum unmolested, with the column down and
the punishments freshly visible. Section 3 introduces a personal
preoccupation that will run through the next several letters:
Cicero wants to slip away to Athens to look in on his student
son, and is suspicious of the tutor Leonides' too-careful
praise (``This is not the testimony of a man who trusts him but
rather of one who fears for him''). Herodes, who was asked to
report \textit{kata miton} --- ``thread by thread'' --- has sent
nothing, which Cicero reads as itself a report.}
